%% =================================================================
%% verzeichnisse/glossar.tex
%% =================================================================

\newglossaryentry{crossplatform}{%
  name={Cross-Platform},
  description={Software-Entwicklungsansatz, bei dem Anwendungen mit einer gemeinsamen Codebasis auf mehreren Betriebssystemen (z. B. Android und iOS) ausführbar sind, hier umgesetzt mit Flutter.}
}

\newglossaryentry{speech2text}{%
  name={Speech-to-Text},
  description={Verfahren zur automatischen Umwandlung gesprochener Sprache in geschriebenen Text, im Projekt realisiert durch das Whisper-Modell.}
}

\newglossaryentry{nlp}{%
  name={Natural Language Processing},
  description={Teilgebiet der künstlichen Intelligenz zur Verarbeitung und Analyse natürlicher Sprache, z. B. für Tokenisierung, POS-Tagging und Informationsextraktion mit spaCy.}
}

\newglossaryentry{bootloader}{%
  name={Bootstrapping / Pipeline-Start},
  description={Initialisierung der Verarbeitungskette im System, bei der Audioaufnahme, Transkription und Analyse sequenziell gestartet werden.}
}

\newglossaryentry{relational}{%
  name={Relationales Modell},
  description={Datenmodell, bei dem Informationen strukturiert in Tabellen gespeichert und über Primär- und Fremdschlüssel (z. B. in SQL-Datenbanken) miteinander verknüpft werden.}
}

\newglossaryentry{offlinefirst}{%
  name={Offline-First},
  description={Architekturansatz, bei dem die Anwendung vollständig lokal funktionsfähig ist und Daten bei bestehender Verbindung synchronisiert werden.}
}

\newglossaryentry{widget}{%
  name={Widget},
  description={Zentrale UI-Baueinheit in Flutter. Jede visuelle oder interaktive Komponente der Streamlit- bzw. Flutter-Oberfläche basiert auf Widgets.}
}

\newglossaryentry{pipeline}{%
  name={KI-Pipeline},
  description={Modulare Verarbeitungskette aus Speech-to-Text, NLP-Analyse und visueller Ausgabe zur schrittweisen Transformation von Sprache in strukturierte Informationen.}
}

% =========================================================
% Dual Entries (Akronyme + Definitionen)
% =========================================================

\newdualentry{api}
{API}
{Application Programming Interface}
{Programmierschnittstelle zur standardisierten Kommunikation zwischen Systemkomponenten, insbesondere für die Anbindung externer Dienste und Bilddatenquellen (z. B. Unsplash API).}

\newdualentry{json}
{JSON}
{JavaScript Object Notation}
{Textbasiertes Datenformat zur strukturierten Übertragung von Informationen zwischen Client (Streamlit-App) und externen oder lokalen Services.}

\newdualentry{stt}
{STT}
{Speech-to-Text}
{Automatische Umwandlung von Audioeingaben in Textdaten mittels Whisper-Modell als erster Schritt der KI-Pipeline.}